\documentclass[10pt,journal]{IEEEtran}
% _preamble.tex --- reglages communs IEEE pour tous les TP du parcours "Improve"
% A inclure via % _preamble.tex --- reglages communs IEEE pour tous les TP du parcours "Improve"
% A inclure via % _preamble.tex --- reglages communs IEEE pour tous les TP du parcours "Improve"
% A inclure via \input{_preamble} JUSTE APRES \documentclass[10pt,journal]{IEEEtran}
% Compilation : tectonic (telecharge les packages automatiquement).
\usepackage[T1]{fontenc}
\usepackage[utf8]{inputenc}
\usepackage[french]{babel}
\usepackage{amsmath,amssymb}
\usepackage{newtxtext,newtxmath}     % rendu Times (proche de Tinos) ; APRES amssymb
\usepackage{graphicx}
\usepackage{booktabs}
\usepackage{array}
\usepackage{xcolor}
\usepackage{listings}
\usepackage{enumitem}
\usepackage{tikz}
\usepackage{pgfplots}
\pgfplotsset{compat=1.18}
\usetikzlibrary{arrows.meta,patterns,calc,decorations.pathmorphing}
\usepackage{cite}
\usepackage[hidelinks]{hyperref}

% --- Noms francais des flottants ---
\renewcommand{\tablename}{Tableau}
\renewcommand{\figurename}{Fig.}
\renewcommand{\lstlistingname}{Listing}
\addto\captionsfrench{\renewcommand{\refname}{Références}}

% --- En-tete de revue commun a tous les TP ---
\newcommand{\improvejournal}{IEEE Transactions on Computational Mechanics and Machine Learning, Vol.~1, No.~1, Juin~2026}

% --- Style code Python (mono, sobre facon IEEE) ---
\definecolor{kw}{HTML}{0033B3}
\definecolor{cm}{HTML}{7A7A7A}
\definecolor{st}{HTML}{067D17}
\lstdefinestyle{py}{%
  language=Python, basicstyle=\ttfamily\footnotesize,
  keywordstyle=\color{kw}, commentstyle=\color{cm}\itshape, stringstyle=\color{st},
  numbers=none, showstringspaces=false, breaklines=true, columns=fullflexible,
  aboveskip=3pt, belowskip=1pt, xleftmargin=2pt,
  literate={é}{{\'e}}1 {è}{{\`e}}1 {à}{{\`a}}1 {ê}{{\^e}}1 {ô}{{\^o}}1 {î}{{\^i}}1 {ç}{{\c c}}1 {É}{{\'E}}1 {_}{{\_}}1}
\lstset{style=py}
 JUSTE APRES \documentclass[10pt,journal]{IEEEtran}
% Compilation : tectonic (telecharge les packages automatiquement).
\usepackage[T1]{fontenc}
\usepackage[utf8]{inputenc}
\usepackage[french]{babel}
\usepackage{amsmath,amssymb}
\usepackage{newtxtext,newtxmath}     % rendu Times (proche de Tinos) ; APRES amssymb
\usepackage{graphicx}
\usepackage{booktabs}
\usepackage{array}
\usepackage{xcolor}
\usepackage{listings}
\usepackage{enumitem}
\usepackage{tikz}
\usepackage{pgfplots}
\pgfplotsset{compat=1.18}
\usetikzlibrary{arrows.meta,patterns,calc,decorations.pathmorphing}
\usepackage{cite}
\usepackage[hidelinks]{hyperref}

% --- Noms francais des flottants ---
\renewcommand{\tablename}{Tableau}
\renewcommand{\figurename}{Fig.}
\renewcommand{\lstlistingname}{Listing}
\addto\captionsfrench{\renewcommand{\refname}{Références}}

% --- En-tete de revue commun a tous les TP ---
\newcommand{\improvejournal}{IEEE Transactions on Computational Mechanics and Machine Learning, Vol.~1, No.~1, Juin~2026}

% --- Style code Python (mono, sobre facon IEEE) ---
\definecolor{kw}{HTML}{0033B3}
\definecolor{cm}{HTML}{7A7A7A}
\definecolor{st}{HTML}{067D17}
\lstdefinestyle{py}{%
  language=Python, basicstyle=\ttfamily\footnotesize,
  keywordstyle=\color{kw}, commentstyle=\color{cm}\itshape, stringstyle=\color{st},
  numbers=none, showstringspaces=false, breaklines=true, columns=fullflexible,
  aboveskip=3pt, belowskip=1pt, xleftmargin=2pt,
  literate={é}{{\'e}}1 {è}{{\`e}}1 {à}{{\`a}}1 {ê}{{\^e}}1 {ô}{{\^o}}1 {î}{{\^i}}1 {ç}{{\c c}}1 {É}{{\'E}}1 {_}{{\_}}1}
\lstset{style=py}
 JUSTE APRES \documentclass[10pt,journal]{IEEEtran}
% Compilation : tectonic (telecharge les packages automatiquement).
\usepackage[T1]{fontenc}
\usepackage[utf8]{inputenc}
\usepackage[french]{babel}
\usepackage{amsmath,amssymb}
\usepackage{newtxtext,newtxmath}     % rendu Times (proche de Tinos) ; APRES amssymb
\usepackage{graphicx}
\usepackage{booktabs}
\usepackage{array}
\usepackage{xcolor}
\usepackage{listings}
\usepackage{enumitem}
\usepackage{tikz}
\usepackage{pgfplots}
\pgfplotsset{compat=1.18}
\usetikzlibrary{arrows.meta,patterns,calc,decorations.pathmorphing}
\usepackage{cite}
\usepackage[hidelinks]{hyperref}

% --- Noms francais des flottants ---
\renewcommand{\tablename}{Tableau}
\renewcommand{\figurename}{Fig.}
\renewcommand{\lstlistingname}{Listing}
\addto\captionsfrench{\renewcommand{\refname}{Références}}

% --- En-tete de revue commun a tous les TP ---
\newcommand{\improvejournal}{IEEE Transactions on Computational Mechanics and Machine Learning, Vol.~1, No.~1, Juin~2026}

% --- Style code Python (mono, sobre facon IEEE) ---
\definecolor{kw}{HTML}{0033B3}
\definecolor{cm}{HTML}{7A7A7A}
\definecolor{st}{HTML}{067D17}
\lstdefinestyle{py}{%
  language=Python, basicstyle=\ttfamily\footnotesize,
  keywordstyle=\color{kw}, commentstyle=\color{cm}\itshape, stringstyle=\color{st},
  numbers=none, showstringspaces=false, breaklines=true, columns=fullflexible,
  aboveskip=3pt, belowskip=1pt, xleftmargin=2pt,
  literate={é}{{\'e}}1 {è}{{\`e}}1 {à}{{\`a}}1 {ê}{{\^e}}1 {ô}{{\^o}}1 {î}{{\^i}}1 {ç}{{\c c}}1 {É}{{\'E}}1 {_}{{\_}}1}
\lstset{style=py}


\begin{document}

\title{Éditorial --- Guide du parcours de travaux pratiques : simulation accélérée par apprentissage automatique}

\author{Prénom~Nom%
\thanks{Profil hybride mécanique + IA --- Parcours Simulation + IA. Document de cadrage (syllabus). Courriel~: prenom.nom@exemple.org}}

\markboth{\improvejournal}{Nom \MakeLowercase{\textit{et~al.}}: Guide du parcours de travaux pratiques}

\maketitle

\begin{abstract}
Ce document de cadrage présente la logique et l'organisation du parcours de travaux pratiques (TP) menant au profil visé : modéliser la physique, l'accélérer par apprentissage automatique, et savoir quand le modèle décroche du réel --- le c\oe{}ur du métier chez un acteur de la simulation tel que CAE. Le parcours se structure en quatre étages de simulation de complexité croissante, traversés par un fil unique : la simulation accélérée par modèle de substitution et sa validation physique. Chaque TP est autonome, se traite au fur et à mesure et produit un livrable concret.
\end{abstract}

\begin{IEEEkeywords}
Syllabus, éléments finis, CFD, dynamique du vol, jumeau numérique, modèle de substitution, validation physique, parcours de formation.
\end{IEEEkeywords}

\IEEEpeerreviewmaketitle

\section{Fil rouge}
\IEEEPARstart{U}{n} seul fil traverse les quatre étages : \emph{la simulation accélérée par apprentissage automatique et sa validation physique}. L'objectif n'est pas de couvrir quatre sujets séparés, mais de bâtir \emph{une} compétence, prouvée quatre fois, sur des domaines physiques de plus en plus larges (Tableau~\ref{tab:etages}).

\begin{table}[!t]
\caption{Les quatre étages de simulation}
\label{tab:etages}
\centering
\begin{tabular}{@{}cll@{}}
\toprule
Étage & Domaine & Aboutissement (portfolio) \\
\midrule
1 & Structures / EF & Substitut d'un solveur EF + validation \\
2 & CFD / aérodynamique & Substitut remplaçant des heures de calcul \\
3 & Dynamique du vol & Modèle de comportement + détection d'écart \\
4 & Jumeau numérique & Physique + données capteurs + ML \\
\bottomrule
\end{tabular}
\end{table}

\section{Étage 1 --- Structures / éléments finis}
On débute ici car c'est le prolongement direct de la conception mécanique, et cela bâtit le squelette réutilisé partout ensuite : discrétisation, maillage, solveur, convergence, validation analytique. La liste des TP figure au Tableau~\ref{tab:tp1}.

\begin{table}[!t]
\caption{Travaux pratiques de l'étage 1}
\label{tab:tp1}
\centering
\begin{tabular}{@{}cll@{}}
\toprule
TP & Titre & Compétence-clé \\
\midrule
1.1 & Flexion : vérité-terrain analytique & Domaine de validité d'un modèle \\
1.2 & Solveur EF 1D (barre) & Matrice de rigidité, assemblage \\
1.3 & Solveur EF poutre + convergence & Maillage, convergence \\
1.4 & Du solveur au substitut & Régression, évaluation honnête \\
1.5 & Détection du hors-domaine & Garde-fou physique \\
\bottomrule
\end{tabular}
\end{table}

Les TP des étages 2 à 4 sont produits au fur et à mesure de la progression.

\section{Méthode de travail}
Chaque TP suit la même structure, pensée pour ancrer le réflexe d'ingénieur : (i) rappel théorique et ses hypothèses (donc ses limites) ; (ii) calcul à la main avant tout code ; (iii) script Python à compléter ; (iv) validation contre la vérité-terrain ; (v) le pont vers l'apprentissage automatique. La tentation de coder immédiatement doit être combattue : la valeur du profil tient à la capacité de \emph{juger} un résultat, ce qui suppose d'avoir fait le calcul à la main au moins une fois. Le code sans la physique est précisément le piège que l'automatisation tend aux ingénieurs.

\section{Critères de réussite}
À la fin de chaque TP, l'apprenant doit pouvoir : énoncer les hypothèses du modèle et son domaine de validité ; retrouver le résultat analytique à la main ; reproduire ce résultat par le code à mieux que 1~\% ; expliquer l'origine de tout écart ; et relier l'exercice à l'étape de substitution et de critique de l'IA.

\section{Conclusion}
Au terme des quatre étages, l'apprenant ne maîtrise pas \og quatre outils\fg{} mais une compétence transférable unique --- la simulation accélérée par ML et sa validation physique --- qui culmine sur le profil recherché en industrie de la simulation. Le prochain document est le TP~1.1 : \emph{Flexion d'une poutre encastrée}.

\end{document}
